\textbf{Phần 1. Trắc nghiệm lựa chọn}
\setcounterpageref{ex}{0}

\begin{ex} %[3D1V05]
Tính giá trị của biểu thức: $A = 3 \times 9 - 27 : 3 + 15$.
\choice
{$30$}
{$36$}
{\True $33$}
{$32$}
\loigiai{
Thực hiện các phép tính nhân và chia trước, cộng trừ sau:\\
$$A = 27 - 9 + 15 = 18 + 15 = 33$$
\textbf{Đáp án của câu hỏi là C}
}
\end{ex}

\begin{ex} %[3D1V05]
Điền dấu thích hợp vào chỗ chấm: $3 \times 7 + 9 \dots 30 : 3 \times 3$.
\choice
{$>$}
{\True $=$}
{$<$}
{Không so sánh được}
\loigiai{
Tính giá trị của hai vế:\\
- Vế trái: $3 \times 7 + 9 = 21 + 9 = 30$\\
- Vế phải: $30 : 3 \times 3 = 10 \times 3 = 30$\\
Vì $30 = 30$ nên ta điền dấu $=$.\\
\textbf{Đáp án của câu hỏi là B}
}
\end{ex}

\begin{ex} %[3D1V05]
Tìm giá trị của $x$, biết: $x : 3 + 18 = 25$.
\choice
{$24$}
{$18$}
{$27$}
{\True $21$}
\loigiai{
Xem $x : 3$ là số hạng chưa biết:\\
$$x : 3 = 25 - 18$$
$$x : 3 = 7$$
$$x = 7 \times 3 = 21$$
\textbf{Đáp án của câu hỏi là D}
}
\end{ex}

\begin{ex} %[3D1V05]
Lớp 3B trồng được $8$ cây hoa. Lớp 3A trồng được số cây hoa gấp $3$ lần lớp 3B. Hỏi cả hai lớp trồng được tất cả bao nhiêu cây hoa?
\choice
{\True $32$ cây}
{$24$ cây}
{$30$ cây}
{$28$ cây}
\loigiai{
Lớp 3A trồng được số cây hoa là:
$$8 \times 3 = 24\text{ (cây)}$$
Cả hai lớp trồng được tất cả số cây hoa là:
$$8 + 24 = 32\text{ (cây)}$$
\textbf{Đáp án của câu hỏi là A}
}
\end{ex}

\begin{ex} %[3D1V05]
Có $27\text{ lít}$ mật ong chia đều vào $3$ can lớn. Người ta lấy $1$ can lớn đó rót đều vào các chai nhỏ, mỗi chai chứa $3\text{ lít}$. Hỏi rót được bao nhiêu chai nhỏ?
\choice
{$4$ chai}
{$2$ chai}
{\True $3$ chai}
{$5$ chai}
\loigiai{
Mỗi can lớn chứa số lít mật ong là:
$$27 : 3 = 9\text{ (lít)}$$
Số chai nhỏ rót được là:
$$9 : 3 = 3\text{ (chai)}$$
\textbf{Đáp án của câu hỏi là C}
}
\end{ex}

\begin{ex} %[3D1V05]
Cho dãy số: $3; 6; 9; 12; \dots; 27$. Số hạng đứng chính giữa của dãy số trên là
\choice
{$12$}
{\True $15$}
{$18$}
{$9$}
\loigiai{
Dãy số gồm các bội số của 3 từ $3 \times 1$ đến $3 \times 9$, có tất cả 9 số hạng.\\
Số hạng đứng chính giữa là số hạng thứ năm:\\
$$3 \times 5 = 15$$
\textbf{Đáp án của câu hỏi là B}
}
\end{ex}

\begin{ex} %[3D1V05]
Trong một phép chia cho $3$, thương bằng $7$ và số dư là số dư lớn nhất có thể có. Số bị chia của phép chia đó là
\choice
{$21$}
{$22$}
{$24$}
{\True $23$}
\loigiai{
Vì số chia là 3 nên số dư lớn nhất có thể có là $2$.\\
Số bị chia là:
$$3 \times 7 + 2 = 21 + 2 = 23$$
\textbf{Đáp án của câu hỏi là D}
}
\end{ex}

\begin{ex} %[3D1V05]
Tính nhanh giá trị của biểu thức: $3 \times 2 + 3 \times 3 + 3 \times 5$.
\choice
{\True $30$}
{$27$}
{$33$}
{$36$}
\loigiai{
Áp dụng tính chất phân phối:\\
$$3 \times 2 + 3 \times 3 + 3 \times 5 = 3 \times (2 + 3 + 5) = 3 \times 10 = 30$$
\textbf{Đáp án của câu hỏi là A}
}
\end{ex}

\begin{ex} %[3D1V05]
Mẹ mua $3$ hộp bánh, mỗi hộp có giá $8$ nghìn đồng và mua thêm $1$ gói kẹo giá $6$ nghìn đồng. Mẹ đưa cho cô bán hàng tờ tiền $50$ nghìn đồng. Hỏi cô bán hàng phải trả lại mẹ bao nhiêu nghìn đồng?
\choice
{$15$ nghìn đồng}
{$25$ nghìn đồng}
{\True $20$ nghìn đồng}
{$10$ nghìn đồng}
\loigiai{
Số tiền mẹ mua 3 hộp bánh là:
$$3 \times 8 = 24\text{ (nghìn đồng)}$$
Tổng số tiền mẹ đã mua bánh và kẹo là:
$$24 + 6 = 30\text{ (nghìn đồng)}$$
Cô bán hàng phải trả lại mẹ số tiền là:
$$50 - 30 = 20\text{ (nghìn đồng)}$$
\textbf{Đáp án của câu hỏi là C}
}
\end{ex}

\begin{ex} %[3D1V05]
Tìm số tự nhiên $a$, biết: $3 \times 6 < a < 3 \times 8$ và $a$ chia hết cho $3$.
\choice
{$19$}
{\True $21$}
{$20$}
{$22$}
\loigiai{
Ta có: $18 < a < 24$.\\
Các số tự nhiên nằm giữa 18 và 24 là: $19; 20; 21; 22; 23$.\\
Trong các số trên, số duy nhất chia hết cho 3 là $21$ (vì $21 : 3 = 7$).\\
\textbf{Đáp án của câu hỏi là B}
}
\end{ex}

\begin{ex} %[3D1V05]
Trên một đoạn đường dài $27\text{ m}$, người ta trồng các cây hoa cách đều nhau $3\text{ m}$, ở cả hai đầu đường đều có trồng cây. Hỏi trên đoạn đường đó có tất cả bao nhiêu cây hoa?
\choice
{$8$ cây}
{$9$ cây}
{$11$ cây}
{\True $10$ cây}
\loigiai{
Số khoảng cách giữa các cây là:
$$27 : 3 = 9\text{ (khoảng)}$$
Vì cả hai đầu đường đều trồng cây nên số cây hoa là:
$$9 + 1 = 10\text{ (cây)}$$
\textbf{Đáp án của câu hỏi là D}
}
\end{ex}

\begin{ex} %[3D1V05]
Tìm giá trị của $x$, biết: $(x - 5) \times 3 = 24$.
\choice
{\True $13$}
{$12$}
{$14$}
{$15$}
\loigiai{
Xem $(x - 5)$ là thừa số chưa biết:\\
$$x - 5 = 24 : 3$$
$$x - 5 = 8$$
$$x = 8 + 5 = 13$$
\textbf{Đáp án của câu hỏi là A}
}
\end{ex}

\textbf{Phần 2. Trắc nghiệm đúng sai}
\setcounterpageref{ex}{0}

\begin{ex} %[3D1V05]
Xét tính đúng sai của các nhận định sau về phép nhân và phép chia với $3$:
\choiceTF
{\True Tích của một số lẻ bất kì nhân với $3$ luôn luôn là một số lẻ.}
{Trong một phép chia cho $3$, số dư lớn nhất có thể có là $3$.}
{\True Nếu một số giảm đi $3$ lần rồi lại gấp lên $3$ lần thì được số ban đầu.}
{\True Tổng của ba số tự nhiên liên tiếp luôn luôn là một số chia hết cho $3$.}
\loigiai{
\begin{itemize}
    \item Ý a) Đúng: Tích của hai số lẻ luôn là số lẻ (ví dụ $1 \times 3 = 3$, $3 \times 3 = 9$, $5 \times 3 = 15$).
    \item Ý b) Sai: Số dư phải luôn bé hơn số chia, vì số chia là 3 nên số dư lớn nhất là 2.
    \item Ý c) Đúng: Vì $a : 3 \times 3 = a$.
    \item Ý d) Đúng: Gọi 3 số liên tiếp là $n, n+1, n+2$. Tổng là $3n + 3 = 3 \times (n+1)$, chia hết cho 3.
\end{itemize}
\begin{itemize}
    \item a) \textbf{Đúng}.
    \item b) \textbf{Sai}.
    \item c) \textbf{Đúng}.
    \item d) \textbf{Đúng}.
\end{itemize}
}
\end{ex}

\begin{ex} %[3D1V05]
Một cửa hàng có $3$ thùng dầu, mỗi thùng chứa $9\text{ lít}$ dầu. Cửa hàng đã bán đi $12\text{ lít}$ dầu. Xét tính đúng sai:
\choiceTF
{\True Ban đầu cửa hàng có tất cả $27\text{ lít}$ dầu.}
{\True Sau khi bán, cửa hàng còn lại $15\text{ lít}$ dầu.}
{Số dầu còn lại đủ để rót đều vào 4 can loại $4\text{ lít}$.}
{\True Số dầu còn lại rót vừa đủ vào 5 chai loại $3\text{ lít}$.}
\loigiai{
\begin{itemize}
    \item Ý a) Đúng: Tổng số dầu ban đầu là: $3 \times 9 = 27\text{ (lít)}$.
    \item Ý b) Đúng: Số dầu còn lại là: $27 - 12 = 15\text{ (lít)}$.
    \item Ý c) Sai: 4 can loại 4 lít cần $4 \times 4 = 16\text{ lít}$, mà cửa hàng chỉ còn 15 lít.
    \item Ý d) Đúng: $15 : 3 = 5$ (rót vừa đủ 5 chai loại 3 lít).
\end{itemize}
\begin{itemize}
    \item a) \textbf{Đúng}.
    \item b) \textbf{Đúng}.
    \item c) \textbf{Sai}.
    \item d) \textbf{Đúng}.
\end{itemize}
}
\end{ex}

\textbf{Phần 3. Trả lời ngắn}
\setcounterpageref{ex}{0}

\begin{ex} %[3D1V05]
Tính giá trị của biểu thức: $3 \times 7 + 18 : 3$.
\shortans{27}
\loigiai{
\textbf{Đáp số:} $27$\\
Thực hiện nhân chia trước, cộng sau:\\
$$3 \times 7 + 18 : 3 = 21 + 6 = 27$$
}
\end{ex}

\begin{ex} %[3D1V05]
Một người thợ cưa một khúc gỗ dài $15\text{ m}$ thành các đoạn dài $3\text{ m}$. Biết mỗi lần cưa người đó mất $5$ phút. Hỏi người đó cưa xong khúc gỗ hết tất cả bao nhiêu phút?
\shortans{20}
\loigiai{
\textbf{Đáp số:} $20$\\
Số đoạn gỗ cưa được là: $15 : 3 = 5$ (đoạn).\\
Để cưa thành 5 đoạn thì chỉ cần cưa: $5 - 1 = 4$ (lần cưa).\\
Thời gian cưa xong khúc gỗ là:
$$4 \times 5 = 20\text{ (phút)}$$
}
\end{ex}

\begin{ex} %[3D1V05]
Tìm số bị chia trong một phép chia cho $3$, biết thương bằng $8$ và số dư bằng $1$.
\shortans{25}
\loigiai{
\textbf{Đáp số:} $25$\\
Số bị chia là:
$$3 \times 8 + 1 = 24 + 1 = 25$$
}
\end{ex}

\begin{ex} %[3D1V05]
Nam có $12$ viên bi. Số viên bi của Việt gấp $3$ lần số viên bi của Nam. Hỏi cả hai bạn có tất cả bao nhiêu viên bi?
\shortans{48}
\loigiai{
\textbf{Đáp số:} $48$\\
Số viên bi của Việt là:
$$12 \times 3 = 36\text{ (viên bi)}$$
Cả hai bạn có tất cả số viên bi là:
$$12 + 36 = 48\text{ (viên bi)}$$
}
\end{ex}

\begin{ex} %[3D1V05]
Tìm giá trị của $x$, biết: $3 \times x = 45 - 18$.
\shortans{9}
\loigiai{
\textbf{Đáp số:} $9$\\
Tính vế phải:\\
$$3 \times x = 27$$
$$x = 27 : 3 = 9$$
}
\end{ex}

\begin{ex} %[3D1V05]
Tính nhanh giá trị của biểu thức: $3 \times 1 + 3 \times 2 + 3 \times 3 + 3 \times 4$.
\shortans{30}
\loigiai{
\textbf{Đáp số:} $30$\\
Đặt thừa số chung:\\
$$3 \times (1 + 2 + 3 + 4) = 3 \times 10 = 30$$
}
\end{ex}
